% §11.5 闭形式与恰当微分形式（全微分）

\section{闭形式与恰当微分形式（全微分）}

\begin{definition}\label{def:11.5.1}

设 $U\subset\mathbb{R}^{n}$ 为开集, $\omega$ 为 $U$ 上的 $k$ 次 $C^{r}\ (r\geqslant1)$ 微分形式. 如果 $\mathrm{d}\omega=0$, 则称 $\omega$ 为\textbf{闭形式}; 如果在 $U$ 上存在 $k-1$ 次 $C^{r+1}$ 微分形式 $\theta$, 使得 $\omega=\mathrm{d}\theta$, 则称 $\omega$ 为\textbf{恰当微分形式} (或\textbf{全微分}).

\end{definition}

\begin{theorem}\label{thm:11.5.1}

(恰当微分形式的必要条件)$\omega$ 为开集 $U\subset\mathbb{R}^{n}$ 上的 $k$ 次 $C^{r}\ (r\geqslant1)$ 恰当微分形式, 则 $\omega$ 为闭形式. 但反之不真.

\end{theorem}

\begin{proof}

因为 $\omega$ 为 $k$ 次 $C^{r}$ 恰当微分形式, 故由定义 \ref{def:11.5.1}, 存在 $k-1$ 次 $C^{r+1}$ 微分形式 $\theta$, 使得 $\omega=\mathrm{d}\theta$. 于是, 根据定理 \ref{thm:11.3.1}(3), 有

\[
\mathrm{d}\omega=\mathrm{d}(\mathrm{d}\theta)=0,
\]

即 $\omega$ 为闭形式.

反之不真. 例如: $U=\mathbb{R}^{2}\setminus\{(0,0)\}$, $\omega=\dfrac{-y\mathrm{d}x+x\mathrm{d}y}{x^{2}+y^{2}}$, $(x,y)\in U$. 由于

\[
\begin{aligned}
\mathrm{d}\omega&=\left(\frac{\partial}{\partial x}\!\left(\frac{x}{x^{2}+y^{2}}\right)-\frac{\partial}{\partial y}\!\left(\frac{-y}{x^{2}+y^{2}}\right)\right)\mathrm{d}x\wedge\mathrm{d}y\\
&=\left[\frac{y^{2}-x^{2}}{(x^{2}+y^{2})^{2}}+\frac{x^{2}-y^{2}}{(x^{2}+y^{2})^{2}}\right]\mathrm{d}x\wedge\mathrm{d}y=0,
\end{aligned}
\]

即 $\omega$ 为闭形式. 但 $\omega$ 却不是恰当微分形式. (反证)假设 $\omega$ 为恰当微分形式, 则由下面的定理 \ref{thm:11.5.2}, 有

\[
\int_{\overrightarrow{C}}\omega=0,
\]

其中 $\overrightarrow{C}$ 为逆时针方向的圆 $x^{2}+y^{2}=R^{2}$. 这与

\[
\begin{aligned}
\int_{\overrightarrow{C}}\omega
&=\int_{\overrightarrow{C}}\frac{-y}{x^{2}+y^{2}}\mathrm{d}x+\frac{x}{x^{2}+y^{2}}\mathrm{d}y\\
&=\int_{0}^{2\pi}\left[\frac{-R\sin\varphi}{R^{2}}(-R\sin\varphi)
  +\frac{R\cos\varphi}{R^{2}}(R\cos\varphi)\right]\mathrm{d}\varphi\\
&=\int_{0}^{2\pi}\mathrm{d}\varphi=2\pi
\end{aligned}
\]

相矛盾.

\end{proof}

\begin{remark}\label{rem:11.5.1}

在 $\mathbb{R}_{+}^{2}=\{(x,y)\mid y>0\}$ 上, 令 $\theta=-\arctan\dfrac{x}{y}$;

在 $\mathbb{R}^{2}\setminus\{(x,0)\mid x\geqslant0\}$ 上, 令

\[
\theta=\begin{cases}
-\arctan\dfrac{x}{y},&y>0,\\[4pt]
\dfrac{\pi}{2},&y=0,\ x<0,\\[4pt]
-\arctan\dfrac{x}{y}+\pi,&y<0,
\end{cases}
\]

则 $\omega=\dfrac{-y}{x^{2}+y^{2}}\mathrm{d}x+\dfrac{x}{x^{2}+y^{2}}\mathrm{d}y=\mathrm{d}\theta$ 为恰当微分形式.

建议读者根据上述结果证明 $\omega$ 在 $\mathbb{R}^{2}\setminus\{(0,0)\}$ 上不为恰当微分形式.

\end{remark}

\begin{theorem}\label{thm:11.5.2}

设 $U\subset\mathbb{R}^{n}$ 为开区域, 则以下各结论等价:

\begin{enumerate}[label=(\arabic*),leftmargin=2.6em,itemsep=0.1\baselineskip,topsep=0.1\baselineskip]

\item $\omega$ 为 $U$ 上的一次 $C^{r}\ (r\geqslant1)$ 恰当微分形式;

\item $\displaystyle\int_{\overrightarrow{C}}\omega=0$, 其中 $\overrightarrow{C}$ 为 $U$ 中任何分段 $\dot{C}^{1}$ 定向简单闭曲线;

\item $\displaystyle\int_{\overrightarrow{C}}\omega$ 在 $U$ 上与道路 $\overrightarrow{C}$ 的选取无关, 而只与 $\overrightarrow{C}$ 的起点 $p$ 和终点 $q$ 有关.

\end{enumerate}

\end{theorem}

\begin{proof}

(1)$\Rightarrow$(2) 设闭曲线 $\overrightarrow{C}$ 的定向参数表示为 $\boldsymbol{x}(t)$, $\alpha\leqslant t\leqslant\beta$, $\boldsymbol{x}(\alpha)=\boldsymbol{x}(\beta)$. 因为 $\omega$ 为恰当微分形式, 故 $\omega=\mathrm{d}\theta$, $\theta$ 为 $0$ 次 $C^{2}$ 微分形式 (函数). 于是,

\[
\int_{\overrightarrow{C}}\omega=\int_{\overrightarrow{C}}\mathrm{d}\theta
=\int_{\overrightarrow{C}}\sum_{i=1}^{n}\theta_{x_{i}}'\mathrm{d}x_{i}
=\int_{\alpha}^{\beta}\sum_{i=1}^{n}\theta_{x_{i}}'\frac{\mathrm{d}x_{i}}{\mathrm{d}t}\mathrm{d}t
=\int_{\alpha}^{\beta}\frac{\mathrm{d}\theta}{\mathrm{d}t}\mathrm{d}t
=\theta(\boldsymbol{x}(\beta))-\theta(\boldsymbol{x}(\alpha))=0.
\]

(2)$\Rightarrow$(3) 设 $\overrightarrow{C}_{1}$ 与 $\overrightarrow{C}_{2}$ 为从 $p$ 到 $q$ 的两条定向 $C^{1}$ 曲线, 则 $\overrightarrow{C}_{1}+\overrightarrow{C}_{2}^{-}$ 为定向闭曲线, 由

\[
0\reason{$(2)$}\int_{\overrightarrow{C}_{1}+\overrightarrow{C}_{2}^{-}}\omega
=\int_{\overrightarrow{C}_{1}}\omega+\int_{\overrightarrow{C}_{2}^{-}}\omega
=\int_{\overrightarrow{C}_{1}}\omega-\int_{\overrightarrow{C}_{2}}\omega
\]

得到

\[
\int_{\overrightarrow{C}_{1}}\omega=\int_{\overrightarrow{C}_{2}}\omega.
\]

(3)$\Rightarrow$(1) 由 (3), 积分 $\int_{\overrightarrow{C}}\omega$ 只与 $\overrightarrow{C}$ 的起点 $p$ 和终点 $q$ 有关, 而与道路 $\overrightarrow{C}$ 的选取无关. 固定起点 $p$, 就得到 $U$ 上的函数

\[
\theta(\boldsymbol{q})=\int_{p}^{q}\omega,\quad\boldsymbol{q}\in U.
\]

则有 $\omega=\mathrm{d}\theta$. 事实上, 设 $\boldsymbol{q}=(x_{1}^{0},x_{2}^{0},\dots,x_{n}^{0})\in U$, $\boldsymbol{q}'=(x_{1}^{0},\dots,x_{i-1}^{0},x_{i}^{0}+\Delta x_{i},x_{i+1}^{0},\dots,x_{n}^{0})$,

\begin{align*}
\frac{\partial\theta}{\partial x_{i}}\bigg|_{q}
&=\lim_{\Delta x_{i}\to0}\frac{1}{\Delta x_{i}}\bigl[\theta(x_{1}^{0},\dots,x_{i-1}^{0},x_{i}^{0}+\Delta x_{i},x_{i+1}^{0},\dots,x_{n}^{0})-\theta(x_{1}^{0},x_{2}^{0},\dots,x_{n}^{0})\bigr]\\
&=\lim_{\Delta x_{i}\to0}\frac{\int_{q}^{q^{\prime}}\omega}{\Delta x_{i}}
 =\lim_{\Delta x_{i}\to0}\frac{\int_{q}^{q^{\prime}}\sum_{i=1}^{n}\omega_{i}\mathrm{d}x_{i}}{\Delta x_{i}}\\
&=\lim_{\Delta x_{i}\to0}\frac{\int_{0}^{1}\omega_{i}(x_{1}^{0},\dots,x_{i-1}^{0},x_{i}^{0}+t\Delta x_{i},x_{i+1}^{0},\dots,x_{n}^{0})\Delta x_{i}\mathrm{d}t}{\Delta x_{i}}\\
&=\omega_{i}(x_{1}^{0},\dots,x_{i-1}^{0},x_{i}^{0},x_{i+1}^{0},\dots,x_{n}^{0}),\quad i=1,2,\dots,n,
\end{align*}

即

\[
\omega=\sum_{i=1}^{n}\omega_{i}\mathrm{d}x_{i}
=\sum_{i=1}^{n}\frac{\partial\theta}{\partial x_{i}}\mathrm{d}x_{i}
=\mathrm{d}\theta
\]

为恰当微分形式.

\end{proof}

\begin{corollary}\label{cor:11.5.1}

设 $U\subsetneqq\mathbb{R}^{n}$ 为开区域, $\omega$ 为 $U$ 上的 $C^{1}$ 恰当微分形式. 如果存在 $0$ 次微分形式 $\theta$, 使得 $\mathrm{d}\theta=\omega$, 则

\[
\theta(\boldsymbol{q})=\int_{p}^{q}\omega+c,\quad\forall\boldsymbol{q}\in U,
\]

其中 $p\in U$ 为固定点, $c$ 为实常数.

\end{corollary}

\begin{proof}

因为

\[
\sum_{i=1}^{n}\frac{\partial\theta}{\partial x_{i}}\mathrm{d}x_{i}
=\mathrm{d}\theta=\omega
=\sum_{i=1}^{n}\omega_{i}\mathrm{d}x_{i},
\]

所以 $\dfrac{\partial\theta}{\partial x_{i}}=\omega_{i}$. 于是

\[
\frac{\partial}{\partial x_{i}}\left(\theta-\int_{p}^{q}\omega\right)
=\frac{\partial\theta}{\partial x_{i}}-\omega_{i}=0.
\]

根据定理 8.1.6 知,

\[
\theta-\int_{p}^{q}\omega=c,\quad\theta=\int_{p}^{q}\omega+c,
\]

其中 $c$ 为实常数.

\end{proof}

\begin{theorem}\label{thm:11.5.3}

(平面 $1$ 次闭形式为恰当微分形式的充分条件)设 $U\subset\mathbb{R}^{2}$ 为单连通 (即 $U$ 中任一条封闭 $C^{1}$ 的 Jordan 曲线内部全在 $U$ 之中, 或任一封闭 $C^{1}$ 的 Jordan 曲线可在 $U$ 中缩成一点)开区域, $\omega$ 为 $U$ 上的一次 $C^{r}\ (r\geqslant1)$ 微分形式. 如果 $\omega$ 为闭形式, 即 $\mathrm{d}\omega=0$, 则 $\omega$ 为恰当微分形式.

\end{theorem}

\begin{proof}

设 $\overrightarrow{C}$ 为 $U$ 中任一封闭 $C^{1}$ 定向曲线, 因 $U$ 是单连通的, 所以 $\overrightarrow{C}$ 围成定向闭区域 $\overrightarrow{D}$, $D\subset U$. 由 Green 公式, 有

\[
\int_{\overrightarrow{C}}\omega=\int_{\overrightarrow{\partial D}}\omega
=\iint_{\overrightarrow{D}}\mathrm{d}\omega
=\iint_{\overrightarrow{D}}0=0.
\]

再由定理 \ref{thm:11.5.2}(2), $\omega$ 为恰当微分形式.

\end{proof}

\begin{remark}\label{rem:11.5.2}

在定理 \ref{thm:11.5.1} 中, $\omega=\dfrac{-y\mathrm{d}x+x\mathrm{d}y}{x^{2}+y^{2}}$, $\mathrm{d}\omega=0$, 即 $\omega$ 为闭形式, 但 $\omega$ 不为恰当微分形式, 这是因为 $U=\mathbb{R}^{2}\setminus\{(0,0)\}$ 不是单连通的开区域.

\end{remark}

\begin{theorem}\label{thm:11.5.4}

($\mathbb{R}^{n}$ 中一次闭形式为恰当微分形式的充分条件)设 $U\subset\mathbb{R}^{n}$ 为一开区域, 且 $U$ 中任一条封闭 $C^{1}$ 定向曲线 $\overrightarrow{C}$ 均可张成一张全在 $U$ 中的二维定向曲面 $\overrightarrow{M}$, $\overrightarrow{\partial M}=\overrightarrow{C}$. $\omega$ 为 $U$ 上的一次 $C^{r}\ (r\geqslant1)$ 微分形式, 如果 $\omega$ 为闭形式, 即 $\mathrm{d}\omega=0$, 则 $\omega$ 为恰当微分形式.

\end{theorem}

\begin{proof}

\[
\int_{\overrightarrow{C}}\omega=\int_{\overrightarrow{\partial M}}\omega
\reason{Stokes 公式}\iint\limits_{\overrightarrow{M}}\mathrm{d}\omega
=\iint\limits_{\overrightarrow{M}}0=0,
\]

再由定理 \ref{thm:11.5.2}(2) 知, $\omega$ 为恰当微分形式.

\end{proof}

\begin{theorem}\label{thm:11.5.5}

($\mathbb{R}^{n}$ 中 $1$ 次闭形式为恰当微分形式的充分条件)设 $U\subset\mathbb{R}^{n}$ 为开区域, 且 $\forall p,q\in U$, 任何连接 $p$ 与 $q$ 的平行于坐标轴的一固定顺序的折线 (例如, 先平行 $x_{1}$ 轴, 再平行 $x_{2}$ 轴, $\dots$, 图 \ref{fig:11.5.1}) 全在 $U$ 中, $\omega$ 为 $U$ 上的一次 $C^{1}$ 闭形式, 即 $\mathrm{d}\omega=0$. 则 $\omega$ 为恰当微分形式, 即 $\omega=\mathrm{d}\theta$.

\end{theorem}

\begin{figure}[htbp]
  \centering
  \begin{tikzpicture}[scale=1.1, >=stealth, line cap=round, line join=round]
    % 区域 U 轮廓
    \draw[thick, fill=ThemeColor!6] plot[smooth cycle, tension=0.85]
      coordinates {(0.3,0.4) (2.2,0.1) (4.2,0.6) (4.5,2.6) (3.2,3.5) (1.2,3.3) (0.2,2.0)};
    \node at (0.8,2.7) {$U$};
    
    % 起点 p, 终点 q
    \coordinate (p) at (1.1,1.0);
    \coordinate (m) at (3.5,1.0);
    \coordinate (q) at (3.5,2.7);
    
    % 折线：先平行 x1 轴，再平行 x2 轴
    \draw[very thick, postaction={decorate}, decoration={markings, mark=at position 0.55 with {\arrow[very thick, ThemeColor]{>}}}]
      (p) -- (m);
    \draw[very thick, postaction={decorate}, decoration={markings, mark=at position 0.55 with {\arrow[very thick, ThemeColor]{>}}}]
      (m) -- (q);
      
    \fill (p) circle (2pt) node[below left] {$p$};
    \fill (q) circle (2pt) node[above right] {$q$};
    \node[below, ThemeColor] at (2.3,1.0) {先平行 $x_{1}$ 轴};
    \node[right, ThemeColor] at (3.5,1.85) {再平行 $x_{2}$ 轴};
  \end{tikzpicture}
  \caption{}\label{fig:11.5.1}
\end{figure}

\begin{proof}

具体构造 $\theta$. 设 $\boldsymbol{x}^{0}=(x_{1}^{0},x_{2}^{0},\dots,x_{n}^{0})\in U$ 为固定点, 令

\[
\theta=\sum_{i=1}^{n}\int_{x_{i}^{0}}^{x_{i}}\omega_{i}(x_{1},x_{2},\dots,x_{i},x_{i+1}^{0},\dots,x_{n}^{0})\,\mathrm{d}x_{i},
\]

则

\[
\begin{aligned}
\frac{\partial\theta}{\partial x_{j}}
&=\frac{\partial}{\partial x_{j}}\Bigl(\sum_{i=1}^{n}\int_{x_{i}^{0}}^{x_{i}}\omega_{i}(x_{1},x_{2},\dots,x_{i},x_{i+1}^{0},\dots,x_{n}^{0})\,\mathrm{d}x_{i}\Bigr)\\
&=\omega_{j}(x_{1},\dots,x_{j},x_{j+1}^{0},\dots,x_{n}^{0})
  +\sum_{i=j+1}^{n}\int_{x_{i}^{0}}^{x_{i}}\frac{\partial\omega_{i}(x_{1},\dots,x_{i},x_{i+1}^{0},\dots,x_{n}^{0})}{\partial x_{j}}\,\mathrm{d}x_{i}\\
&\reason{$\mathrm{d}\omega=0\Leftrightarrow\frac{\partial\omega_{i}}{\partial x_{j}}=\frac{\partial\omega_{j}}{\partial x_{i}}$}\omega_{j}(x_{1},\dots,x_{j},x_{j+1}^{0},\dots,x_{n}^{0})\\
&\quad+\sum_{i=j+1}^{n}\int_{x_{i}^{0}}^{x_{i}}\frac{\partial\omega_{j}(x_{1},\dots,x_{i},x_{i+1}^{0},\dots,x_{n}^{0})}{\partial x_{i}}\,\mathrm{d}x_{i}\\
&=\omega_{j}(x_{1},\dots,x_{j},x_{j+1}^{0},\dots,x_{n}^{0})\\
&\quad+\sum_{i=j+1}^{n}\bigl[\omega_{j}(x_{1},\dots,x_{i},x_{i+1}^{0},\dots,x_{n}^{0})\\
&\qquad-\omega_{j}(x_{1},\dots,x_{i-1},x_{i}^{0},x_{i+1}^{0},\dots,x_{n}^{0})\bigr]\\
&=\omega_{j}(x_{1},\dots,x_{n}).
\end{aligned}
\]

于是,

\[
\mathrm{d}\theta=\sum_{j=1}^{n}\frac{\partial\theta}{\partial x_{j}}\mathrm{d}x_{j}
=\sum_{j=1}^{n}\omega_{j}\mathrm{d}x_{j}=\omega.\qedhere
\]

\end{proof}

\begin{theorem}\label{thm:11.5.6}

设 $U\subset\mathbb{R}^{3}$ 为开区域, $(x_{0},y_{0},z_{0})\in U$ 为一固定点, $\forall(x,y,z)\in U$, 连接 $(x_{0},y_{0},z_{0})$ 与 $(x_{0},y,z_{0})$, 连接 $(x_{0},y,z_{0})$ 与 $(x,y,z_{0})$, 连接 $(x,y,z_{0})$ 与 $(x,y,z)$ 得到的折线全在 $U$ 中 (如 $U$ 为开长方体, 第一卦限, 半空间, 图 \ref{fig:11.5.2}). 如果

\[
\omega=P\mathrm{d}y\wedge\mathrm{d}z+Q\mathrm{d}z\wedge\mathrm{d}x+R\mathrm{d}x\wedge\mathrm{d}y
\]

为 $U$ 上的 $2$ 次 $C^{r}$ 闭微分形式 $(r\geqslant1)$, 即 $\mathrm{d}\omega=0$, 则 $\omega$ 为恰当微分形式, 即 $\omega=\mathrm{d}\theta$.

\end{theorem}

\begin{figure}[htbp]
  \centering
  \begin{tikzpicture}[scale=1.15, >=stealth, line cap=round, line join=round]
    % 长方体 U 顶点
    \coordinate (O) at (0,0);
    \coordinate (D) at (2.8,0);
    \coordinate (B) at (3.6,1.4);
    \coordinate (A) at (0.8,1.4);
    
    \coordinate (E) at (0,1.9);
    \coordinate (G) at (2.8,1.9);
    \coordinate (C) at (3.6,3.3);
    \coordinate (F) at (0.8,3.3);
    
    % 长方体浅色面与边界
    \fill[ThemeColor!4] (O)--(D)--(G)--(E)--cycle;
    \fill[ThemeColor!8] (E)--(G)--(C)--(F)--cycle;
    \fill[ThemeColor!12] (D)--(B)--(C)--(G)--cycle;
    
    \draw[thin, gray!70] (O)--(D)--(B)--(A)--cycle;
    \draw[thin, gray!70] (E)--(G)--(C)--(F)--cycle;
    \draw[thin, gray!70] (O)--(E) (D)--(G) (B)--(C) (A)--(F);
    
    % 三段折线：(x0,y0,z0) -> (x,y0,z0) -> (x,y,z0) -> (x,y,z)
    \coordinate (P0) at (O);
    \coordinate (P1) at (D);
    \coordinate (P2) at (B);
    \coordinate (P3) at (C);
    
    \draw[very thick, postaction={decorate}, decoration={markings, mark=at position 0.55 with {\arrow[very thick, ThemeColor]{>}}}]
      (P0) -- (P1);
    \draw[very thick, postaction={decorate}, decoration={markings, mark=at position 0.55 with {\arrow[very thick, ThemeColor]{>}}}]
      (P1) -- (P2);
    \draw[very thick, postaction={decorate}, decoration={markings, mark=at position 0.55 with {\arrow[very thick, ThemeColor]{>}}}]
      (P2) -- (P3);
      
    \fill (P0) circle (2pt) node[below left] {$(x_{0},y_{0},z_{0})$};
    \fill (P3) circle (2pt) node[above right] {$(x,y,z)$};
    \node at (0.5,2.4) {$U$};
  \end{tikzpicture}
  \caption{}\label{fig:11.5.2}
\end{figure}

\begin{proof}

我们只需求出一个 $\theta$, 使 $\omega=\mathrm{d}\theta$. 当然应先到简单的形式去找. 为此, 试一试能否求出形如

\[
\theta=u\mathrm{d}x+v\mathrm{d}y\quad(\mathrm{d}z\text{ 的系数为 }0)
\]

的一次 $C^{2}$ 微分形式, 使得

\[
\begin{aligned}
&P\mathrm{d}y\wedge\mathrm{d}z+Q\mathrm{d}z\wedge\mathrm{d}x+R\mathrm{d}x\wedge\mathrm{d}y\\
&\quad=\omega=\mathrm{d}(u\mathrm{d}x+v\mathrm{d}y)\\
&\quad=\left(\frac{\partial u}{\partial y}\mathrm{d}y+\frac{\partial u}{\partial z}\mathrm{d}z\right)\wedge\mathrm{d}x
   +\left(\frac{\partial v}{\partial x}\mathrm{d}x+\frac{\partial v}{\partial z}\mathrm{d}z\right)\wedge\mathrm{d}y\\
&\quad=-\frac{\partial v}{\partial z}\mathrm{d}y\wedge\mathrm{d}z
     +\frac{\partial u}{\partial z}\mathrm{d}z\wedge\mathrm{d}x
     +\left(\frac{\partial v}{\partial x}-\frac{\partial u}{\partial y}\right)\mathrm{d}x\wedge\mathrm{d}y.
\end{aligned}
\tag{1}
\]

由此得

\[
\begin{cases}
-\dfrac{\partial v}{\partial z}=P,\\[6pt]
\dfrac{\partial u}{\partial z}=Q,\\[6pt]
\dfrac{\partial v}{\partial x}-\dfrac{\partial u}{\partial y}=R.
\end{cases}
\tag{2}
\]

即

\[
R(x,y,z)=\frac{\partial v}{\partial x}-\frac{\partial u}{\partial y}.
\tag{3}
\]

设 $(x_{0},y_{0},z_{0})\in U$ 为固定点, $\forall(x,y,z)\in U$. 由题设和式 (1), (2) 可取

\[
v(x,y,z)=-\int_{z_{0}}^{z}P(x,y,z)\,\mathrm{d}z,
\tag{4}
\]

\[
u(x,y,z)=\int_{z_{0}}^{z}Q(x,y,z)\,\mathrm{d}z+f(x,y),
\tag{5}
\]

其中 $f$ 为 $U$ 上任意 $C^{1}$ 函数. 将式 (4), (5) 代入式 (3), 并由

\[
\begin{aligned}
0&=\mathrm{d}\omega=\left(\frac{\partial P}{\partial x}+\frac{\partial Q}{\partial y}+\frac{\partial R}{\partial z}\right)\mathrm{d}x\wedge\mathrm{d}y\wedge\mathrm{d}z\\
&\Leftrightarrow\quad\frac{\partial P}{\partial x}+\frac{\partial Q}{\partial y}=-\frac{\partial R}{\partial z}
\end{aligned}
\tag{6}
\]

得到

\[
\begin{aligned}
R(x,y,z)&\reason{$(3)$}\frac{\partial v}{\partial x}-\frac{\partial u}{\partial y}
\reason{$(4),(5)$}-\int_{z_{0}}^{z}\left(\frac{\partial P}{\partial x}+\frac{\partial Q}{\partial y}\right)\mathrm{d}z-\frac{\partial f}{\partial y}
\reason{$(6)$}\int_{z_{0}}^{z}\frac{\partial R}{\partial z}\,\mathrm{d}z-\frac{\partial f}{\partial y}\\
&=R(x,y,z)-R(x,y,z_{0})-\frac{\partial f}{\partial y},
\end{aligned}
\]

所以,

\[
\frac{\partial f}{\partial y}(x,y)=-R(x,y,z_{0}).
\]

取

\[
f(x,y)=-\int_{y_{0}}^{y}R(x,y,z_{0})\,\mathrm{d}y,
\]

再将此式代入 (5) 得到

\[
u(x,y,z)=\int_{z_{0}}^{z}Q(x,y,z)\,\mathrm{d}z-\int_{y_{0}}^{y}R(x,y,z_{0})\,\mathrm{d}y.
\]

验证上述 $u,v$ 满足式 (1), (2); (3), 即得 $\mathrm{d}\theta=\omega$.

\end{proof}

\begin{example}\label{ex:11.5.1}

证明: $\omega=(x+y+1)\mathrm{d}x+(x-y^{2}+3)\mathrm{d}y$ 为恰当微分形式; 并求出 $\theta$, 使 $\mathrm{d}\theta=\omega$. 进而求全 (恰当)微分方程 $(x+y+1)\mathrm{d}x+(x-y^{2}+3)\mathrm{d}y=0$ 的解.

\end{example}

\begin{solution}

设 $P=x+y+1$, $Q=x-y^{2}+3$, 则

\[
\frac{\partial Q}{\partial x}-\frac{\partial P}{\partial y}=1-1=0,
\]

故 $\mathrm{d}\omega=\mathrm{d}(P\mathrm{d}x+Q\mathrm{d}y)=\left(\frac{\partial Q}{\partial x}-\frac{\partial P}{\partial y}\right)\mathrm{d}x\wedge\mathrm{d}y=0$, 即 $\omega$ 为 $C^{2}$ 闭形式. 又因 $U=\mathbb{R}^{2}$ 为单连通区域, 根据定理 \ref{thm:11.5.3}, $\omega$ 为恰当微分形式.

取固定点 $(0,0)\in\mathbb{R}^{2}$, $\forall(x,y)\in\mathbb{R}^{2}$, 有

\[
\begin{aligned}
\theta(x,y)&=\int_{\overrightarrow{C}}\omega\\
&=\left(\int_{\overrightarrow{C}_{1}}+\int_{\overrightarrow{C}_{2}}\right)(x+y+1)\mathrm{d}x+(x-y^{2}+3)\mathrm{d}y\\
&=\int_{0}^{x}(x+0+1)\,\mathrm{d}x+\int_{0}^{y}(x-y^{2}+3)\,\mathrm{d}y\\
&=\left(\frac{x^{2}}{2}+x\right)+\left(xy-\frac{y^{3}}{3}+3y\right)\\
&=\frac{x^{2}}{2}+x+xy-\frac{y^{3}}{3}+3y.
\end{aligned}
\]

进而, 知

\[
\mathrm{d}\theta=\omega=(x+y+1)\mathrm{d}x+(x-y^{2}+3)\mathrm{d}y=0
\]

的解为

\[
\frac{x^{2}}{2}+x+xy-\frac{y^{3}}{3}+3y=\theta(x,y)=c\quad(\text{常数}),
\]

这是全 (恰当)微分方程 $(x+y+1)\mathrm{d}x+(x-y^{2}+3)\mathrm{d}y=0$ 的解.

\end{solution}

\begin{example}\label{ex:11.5.2}

证明: $\omega=(xy+1)\mathrm{d}y\wedge\mathrm{d}z+z\mathrm{d}z\wedge\mathrm{d}x-yz\mathrm{d}x\wedge\mathrm{d}y$ 为恰当微分形式; 并求出 $\theta$, 使 $\mathrm{d}\theta=\omega$.

\end{example}

\begin{solution}

因为

\[
\mathrm{d}\omega=\left(\frac{\partial P}{\partial x}+\frac{\partial Q}{\partial y}+\frac{\partial R}{\partial z}\right)\mathrm{d}x\wedge\mathrm{d}y\wedge\mathrm{d}z
=(y+0-y)\mathrm{d}x\wedge\mathrm{d}y\wedge\mathrm{d}z=0,
\]

故 $\omega$ 为闭形式. 又因 $U=\mathbb{R}^{3}$, 根据定理 \ref{thm:11.5.6} 知, $\omega$ 为恰当微分形式.

为求 $\theta=u\mathrm{d}x+v\mathrm{d}y$, 套用定理 \ref{thm:11.5.6} 证明中的步骤与方法, 有

\[
\begin{cases}
-\dfrac{\partial v}{\partial z}=xy+1&\Rightarrow v=-\displaystyle\int_{0}^{z}(xy+1)\,\mathrm{d}z=-(xy+1)z,\\[10pt]
\dfrac{\partial u}{\partial z}=z&\Rightarrow u=\displaystyle\int_{0}^{z}z\,\mathrm{d}z+f(x,y)=\frac{z^{2}}{2}+f(x,y),\\[10pt]
\dfrac{\partial v}{\partial x}-\dfrac{\partial u}{\partial y}=-yz&\Rightarrow -yz-\dfrac{\partial f}{\partial y}=-yz,\text{ 即 }\dfrac{\partial f}{\partial y}=0,\text{ 取 }f(x,y)=0.
\end{cases}
\]

于是,

\[
\begin{aligned}
\theta&=u\mathrm{d}x+v\mathrm{d}y=\left[\frac{z^{2}}{2}+f(x,y)\right]\mathrm{d}x-(xy+1)z\mathrm{d}y\\
&=\frac{z^{2}}{2}\mathrm{d}x-(xy+1)z\mathrm{d}y.
\end{aligned}
\]

当然, 也可套用定理 \ref{thm:11.5.6} 证明中关于 $u,v$ 的公式, 计算时可快一些. 但是, 没有书时, 记住方法自己当场能推导出来. 各有优缺点.

\end{solution}

\begin{example}\label{ex:11.5.3}

验证第二型曲线积分

\[
\int_{\overrightarrow{C}}(y+z)\mathrm{d}x+(z+x)\mathrm{d}y+(x+y)\mathrm{d}z
\]

与路径无关, 并求被积形式的原函数 $\theta(x,y,z)$.

\end{example}

\begin{solution}

设 $\omega=(y+z)\mathrm{d}x+(z+x)\mathrm{d}y+(x+y)\mathrm{d}z$, 则

\[
\mathrm{d}\omega=(\mathrm{d}y+\mathrm{d}z)\wedge\mathrm{d}x+(\mathrm{d}z+\mathrm{d}x)\wedge\mathrm{d}y+(\mathrm{d}x+\mathrm{d}y)\wedge\mathrm{d}z=0,
\]

即 $\omega$ 为闭形式, 又因为 $U=\mathbb{R}^{3}$ 满足定理 \ref{thm:11.5.5} 中的条件, 故有

\[
\begin{aligned}
\theta(x,y,z)&=\int_{0}^{x}(0+0)\,\mathrm{d}x+\int_{0}^{y}(0+x)\,\mathrm{d}y+\int_{0}^{z}(x+y)\,\mathrm{d}z\\
&=0+xy+(x+y)z=xy+xz+yz.
\end{aligned}
\]

由推论 \ref{cor:11.5.1} 知, 一般的原函数为

\[
\theta(x,y,z)=xy+xz+yz+c,
\]

其中 $c$ 为任意常数 (实际上用观察法也能看出原函数为 $xy+xz+yz+c$).

\end{solution}

\begin{example}\label{ex:11.5.4}

设 $f$ 为 $\mathbb{R}$ 上的 $C^{1}$ 函数, $f(0)=-\dfrac{1}{2}$, 求 $f(x)$ 使第二型曲线积分

\[
\int_{A}^{B}\bigl[\mathrm{e}^{x}+f(x)\bigr]y\mathrm{d}x-f(x)\mathrm{d}y
\]

与路径无关, 只与起点 $A$ 和终点 $B$ 有关. 并求当 $A,B$ 分别为 $(0,0)$, $(1,1)$ 时, 此第二型曲线积分的值.

\end{example}

\begin{solution}

因为要求第二型曲线积分

\[
\int_{A}^{B}\bigl[\mathrm{e}^{x}+f(x)\bigr]y\mathrm{d}x-f(x)\mathrm{d}y
\]

与路径无关, 只与起点 $A$ 和终点 $B$ 有关, 所以有

\[
0=\mathrm{d}\Bigl(\bigl[\mathrm{e}^{x}+f(x)\bigr]y\mathrm{d}x-f(x)\mathrm{d}y\Bigr)
=-\bigl(\mathrm{e}^{x}+f(x)\bigr)\mathrm{d}x\wedge\mathrm{d}y-f'(x)\mathrm{d}x\wedge\mathrm{d}y,
\]

即

\[
f'(x)+f(x)=-\mathrm{e}^{x}.
\]

其解为

\[
f(x)=\mathrm{e}^{-x}\left(\int-\mathrm{e}^{x}\cdot\mathrm{e}^{x}\,\mathrm{d}x+c\right)
=\mathrm{e}^{-x}\left(-\frac{1}{2}\mathrm{e}^{2x}+c\right)
=-\frac{1}{2}\mathrm{e}^{x}+c\mathrm{e}^{-x}.
\]

再由 $f(0)=-\dfrac{1}{2}$ 得到

\[
-\frac{1}{2}=-\frac{1}{2}\mathrm{e}^{0}+c\mathrm{e}^{0}=-\frac{1}{2}+c,\quad c=0,
\]

所以 $f(x)=-\dfrac{1}{2}\mathrm{e}^{x}$.

当 $A,B$ 的坐标分别为 $(0,0)$, $(1,1)$ 时, 有

\[
\begin{aligned}
\int_{(0,0)}^{(1,1)}\left(\mathrm{e}^{x}-\frac{1}{2}\mathrm{e}^{x}\right)y\mathrm{d}x+\frac{1}{2}\mathrm{e}^{x}\mathrm{d}y
&=\int_{(0,0)}^{(1,1)}\frac{1}{2}\mathrm{e}^{x}y\mathrm{d}x+\frac{1}{2}\mathrm{e}^{x}\mathrm{d}y\\
&=\int_{(0,0)}^{(1,0)}0\,\mathrm{d}x+\int_{(1,0)}^{(1,1)}\frac{1}{2}\mathrm{e}^{1}\,\mathrm{d}y
=0+\frac{\mathrm{e}}{2}=\frac{\mathrm{e}}{2}.
\end{aligned}\qedhere
\]

\end{solution}

\begin{exblock}{练习题 11.5}
\item 计算下列恰当微分形式的第二型曲线积分:

  \begin{enumerate}[label=(\arabic*),leftmargin=2.6em,itemsep=0.15\baselineskip,topsep=0.2\baselineskip]

  \item $\displaystyle\int_{(1,1,1)}^{(2,3,-4)}x\mathrm{d}x+y^{2}\mathrm{d}y-z^{2}\mathrm{d}z$;

  \item $\displaystyle\int_{(1,2,3)}^{(0,1,1)}yz\,\mathrm{d}x+xz\,\mathrm{d}y+xy\,\mathrm{d}z$;

  \item $\displaystyle\int_{(x_{1},y_{1},z_{1})}^{(x_{2},y_{2},z_{2})}\frac{x\mathrm{d}x+y\mathrm{d}y+z\mathrm{d}z}{\sqrt{x^{2}+y^{2}+z^{2}}}$, 其中 $(x_{1},y_{1},z_{1})$ 是球面 $x^{2}+y^{2}+z^{2}=a^{2}\ (a>0)$ 上的一点, $(x_{2},y_{2},z_{2})$ 是球面 $x^{2}+y^{2}+z^{2}=b^{2}\ (b>0)$ 上的一点.

  \end{enumerate}

\item 设 $f,g,h$ 为 $\mathbb{R}$ 上的单变量的 $C^{r}\ (r\geqslant1)$ 函数. 计算第二型曲线积分

  \[
  \int_{(x_{1},y_{1},z_{1})}^{(x_{2},y_{2},z_{2})}f(x)\mathrm{d}x+g(y)\mathrm{d}y+h(z)\mathrm{d}z.
  \]

  如果 $f$ 为 $\mathbb{R}$ 上的连续函数, 结论如何?

\item 设 $f$ 为 $\mathbb{R}$ 上的单变量的 $C^{r}\ (r\geqslant1)$ 函数. 计算第二型曲线积分:

  \begin{enumerate}[label=(\arabic*),leftmargin=2.6em,itemsep=0.15\baselineskip,topsep=0.2\baselineskip]

  \item $\displaystyle\int_{(x_{1},y_{1},z_{1})}^{(x_{2},y_{2},z_{2})}f(x+y+z)(\mathrm{d}x+\mathrm{d}y+\mathrm{d}z)$;

  \item $\displaystyle\int_{(x_{1},y_{1},z_{1})}^{(x_{2},y_{2},z_{2})}f\bigl(\sqrt{x^{2}+y^{2}+z^{2}}\bigr)(x\mathrm{d}x+y\mathrm{d}y+z\mathrm{d}z)$.

  \end{enumerate}

  如果 $f$ 为 $\mathbb{R}$ 上的连续函数, 结论如何?

\item 求解下列恰当微分方程:

  \begin{enumerate}[label=(\arabic*),leftmargin=2.6em,itemsep=0.15\baselineskip,topsep=0.2\baselineskip]

  \item $x\mathrm{d}x+y\mathrm{d}y=0$;

  \item $x\mathrm{d}y+y\mathrm{d}x=0$;

  \item $(x+2y)\mathrm{d}x+(2x+y)\mathrm{d}y=0$;

  \item $(x^{2}-y)\mathrm{d}x-(x+\sin^{2}y)\mathrm{d}y=0$;

  \item $\mathrm{e}^{y}\mathrm{d}x+(x\mathrm{e}^{y}-2y)\mathrm{d}y=0$;

  \item $(x+y+1)\mathrm{d}x+(x-y^{2}+3)\mathrm{d}y=0$;

  \item $\dfrac{x\mathrm{d}x+y\mathrm{d}y}{\sqrt{x^{2}+y^{2}}}=\dfrac{y\mathrm{d}x-x\mathrm{d}y}{x^{2}}$;

  \item $\dfrac{x\mathrm{d}y-y\mathrm{d}x}{x^{2}+y^{2}}=x\mathrm{d}x+y\mathrm{d}y$.

  \end{enumerate}
\end{exblock}
